Для количественной оценки качества распознавания текста использовалась метрика WER. Данная метрика предназначена для сравнения распознанного текста с эталонным вариантом и показывает, насколько результат OCR отличается от правильной текстовой расшифровки на уровне слов~\cite{refOCRMetrics}. WER удобна тем, что учитывает наиболее важные для практической задачи типы ошибок: замену одного слова другим, пропуск слова и вставку лишнего слова~\cite{refOCRMetrics}. Благодаря этому метрика позволяет оценивать не только посимвольные искажения как метрика CER~\cite{refOCRMetrics}, но и итоговое качество текста как целостной последовательности слов, поскольку изменение одной буквы в слове может привести к изменению семантического значения слова.

Преимущество WER состоит в том, что она хорошо отражает пригодность распознанного текста для дальнейшего использования. Даже если отдельные символы распознаны не полностью корректно, итоговое значение метрики показывает, насколько сильно искажено содержание текста в целом. Поэтому WER является удобным инструментом для сравнительного анализа различных конфигураций OCR-конвейера. Следует отдельно отметить, что в классической интерпретации WER меньшие значения соответствуют лучшему качеству распознавания. Далее в работе для численного значения качества обработки OCR используется $1-\mathtt{WER}$, где 1 это полное сходство текстов. 